% =============================================================================
% discussion.tex — Policy Discussion section (draft v0.3, 2026-09-03)
% =============================================================================

\section{Policy Discussion}
\label{sec:discussion}

\subsection{Why a linear carbon tax fails and what works instead}

The baseline system is demand-constrained: at the equilibrium, the binding
constraints are the biochar demand segments, not biomass availability. In this
regime a linear tax $\tau E^{+}(x)$ on gross emissions is a near-constant shift
of the objective (gross emissions change little across feasible dispatches), so
the primal optimum does not move, precisely the failure mode characterized by
\cite{benjaafar2013}. Proposition~\ref{prop:taxinv} states the boundary
condition precisely: the tax is dispatch-neutral up to a threshold whose
denominator is the marginal price gap divided by the path emission intensity,
which we compute at $\approx$\$318/t CO$_2$e ($\approx$\$464/t in the
mature-market scenario) with route-specific marginal
intensities of the LP basis (SI Section~A.11), above every tested rate
(\$25--200/t) but close enough that substantially higher taxes could bite. The
free-facility re-solve (C1-free, Section~\ref{sec:policyC})
confirms that this null response is economic rather than an artifact of the
fixed layout. The claim's boundary condition is equally important: it holds
because demand binds; in a supply-constrained variant of the system
(e.g., feedstock scarcity binding before demand) a tax would bite, and the
statement should not be read as tax-neutrality of biochar systems in general.
Our formulation shows that policy instruments which
scale with throughput (baseline-and-credit) or which act as
constraints (gross emission caps) do change the solution: credits
activate the \$250/t bulk-agricultural demand segment at $p_c=\$25$/tCO$_2$e
under the calibrated step demand and
mobilize essentially all near-term biomass at \$200/tCO$_2$e, while caps reshape
production along the marginal abatement cost curve. The net-flux cap
experiment completes the picture: every tested cap
($\mathrm{CAP}^{\mathrm{net}}\ge-1{,}000$~kt) is looser than the baseline flux
of $-1{,}595$~kt and therefore does not bind; the frontier sweep of SI
Fig.~S2 shows that caps tightened below the baseline flux bind in two
stages---expansion into the L segment at $\approx$\$110/tCO$_2$e, then
500$^\circ$C conversion at $\approx$\$150--330/tCO$_2$e. Policy
design for negative-emission supply chains must therefore
explicitly value removal, not merely penalize emissions.

\subsection{Shadow carbon prices versus real markets}

The conditional fixed-layout shadow marginal abatement cost traced by the
gross-emission cap sweep
is $\approx$ \$351--391/t CO$_2$e over the first binding intervals and rises to
$\approx$ \$599/t as the cap tightens; the cap's vertex dual reproduces it
to numerical tolerance (the reported value is the negative of the marginal
value under
Gurobi's sign convention), validating the shadow-price
interpretation on a full supply chain network. This places the Wisconsin
biochar
system's gross-emission abatement above 2024--2025 RGGI allowance prices
(\$17--25/t), EU ETS prices (\texteuro 50--90/t, $\approx$\$55--100/t), and the
voluntary removal
credit band (\$100--250/tCO$_2$e for VM0044/Puro.earth credits). Two
implications follow. First, at any observed carbon price, a cap-and-trade
program alone would induce little gross abatement in this system, because
gross emissions are intrinsic to drying and pyrolysis and abating them means
forgoing biochar sales; reaching the \$800/t-market equilibrium required no
policy at all, but expanding to the bulk agricultural segment requires credit
prices above roughly \$25/tCO$_2$e under the calibrated demand, well within
observed voluntary-market
prices. Second, the comparison of the system's shadow abatement cost with
administered market prices shows that removal crediting, not
gross-emission pricing, is the mechanism that can profitably expand biochar
CDR, a quantitative argument for designing negative-emission policies around
the removal, rather than around the gross emissions, of such systems.

\subsection{Correcting a common policy misattribution: 45Q}

We note explicitly that the U.S.~45Q tax credit, frequently cited in biochar
investment documents, does not apply to biochar. 45Q credits captured
CO$_2$ streams (point-source CCS and DAC); biochar carbon is sequestration, not
capture, and falls outside the statute. Biochar removal credits are instead
issued under voluntary methodologies (Verra VM0044, which requires molar
H/C~$\le$~0.7 and centennial durability, and Puro.earth CORC); VM0044 credits
were endorsed by the ICVCM as high-integrity in 2025 \cite{icvcm2025}. The
policy landscape is, however, in motion: legislation under consideration in the
119th Congress (e.g., the Wildfire Reduction and Carbon Removal Act of 2025,
S.~1842) would create a federal biochar carbon-removal incentive and reform
45Q, which could materially change the economics quantified here
\cite{wrca2025}. Our Paradigm C2 operationalizes the VM0044 gate with
literature-consistent H/C values: 300$^\circ$C biochar is torrefaction-like
(H/C $\approx$0.9--1.4) and earns no credits for any feedstock, while
all 500$^\circ$C chars (H/C $\approx$0.40--0.52) qualify. The model responds
by converting the system to 500$^\circ$C pyrolysis as credit prices rise: a
partial switch already at \$50/t in the mature-market scenario, complete at
\$100/t, and full conversion (2.25~Mt near-term, 4.25~Mt
mature-market) at \$200/t; the conversion
buys durability at the cost of quantity: near-term production falls slightly
below the 2.30~Mt H+M demand target because 500$^\circ$C yields less char per
dry tonne. The gate switches technology without ever subsidizing
low-durability biochar, an example of
policy detail directly reshaping technology selection.

\subsection{Technology selection and quality differentiation}

Two mechanisms drive pyrolysis temperature choice, and which one binds is
itself the finding. First, when dry feedstock is scarce (near-term), per-dry-tonne
margins rule: 300$^\circ$C dominates for the bulk (L) segment up to
$p_c\approx\$239$/t, because its yield advantage (0.352 vs 0.278~t/t)
outweighs the modest credit-rate premium of 500$^\circ$C biochar. Second, when
feedstock is abundant (mature-market), the binding resource is biochar demand,
and 500$^\circ$C's lower yield becomes an asset: each tonne of biochar
served requires more dry tonnes, each earning credits, so 500$^\circ$C overtakes
300$^\circ$C at $p_c\approx\$61$/t for every segment; the model switches
$\approx$2.8~Mt of production at \$100--\$200/t (the \$100/t point from a
dedicated 3{,}600~s re-solve at a 1.9\% gap, archived as
\texttt{policy\_A\_recert2\_v2.csv}). Third, quality-differentiated
demand (v2b) adds a structural floor: the \$800/t premium segment is BC500-only,
so the no-policy baseline already builds 0.80~Mt/yr of BC500 capacity, and
varying the premium capacity moves BC500 one-for-one (0.40~Mt at half capacity,
1.20~Mt at 1.5$\times$). The market structure (whether buyers differentiate
biochar by durability) therefore matters as much as the carbon price, and the
resource regime decides which margin arithmetic applies.

\subsection{Implications for Wisconsin and the Midwest}

The analysis yields three region-specific messages; magnitudes are calibrated
to Wisconsin, while the mechanisms should transfer qualitatively across the
Upper Midwest (Section~\ref{sec:limitations}). (i) Demand, not feedstock,
is the binding constraint: Wisconsin's near-term biomass could support roughly
double the equilibrium biochar production if the \$250/t bulk market is
activated; policies that subsidize farmer adoption (or that value carbon at
voluntary-market rates) expand the system more effectively than supply-side
measures. (ii) Facility locations concentrate in the corn belt and shift with
supply scenarios: mature-market supply lowers biochar's spatial inherent value
from \$324 to \$275/t, flattening location rents. (iii) The dairy-dominated
waste streams (manure, paper-mill sludge) that characterize Wisconsin are
outside this pyrolysis-focused model but form a natural co-location layer for
anaerobic-digestion integrated bioeconomies, an explicit direction for the
multi-pathway extension.

\subsection{Limitations}
\label{sec:limitations}

Several assumptions bound the interpretation of these results. (i)
\emph{Market clearing.} The model solves a single coordinated system optimum
of Sampat et al.~\cite{sampat2019}; the dual-based ``equilibrium prices'' are marginal
values under perfect coordination, not outcomes of a decentralized Nash market,
so magnitudes (especially policy responses) should be read as
coordination-consistent benchmarks. (ii) \emph{Technology space.} Only slow
pyrolysis at two temperatures (300/500$^\circ$C) is modeled; intermediate
temperatures and competing conversion routes (gasification, hydrothermal
carbonization) are excluded, so the technology-choice conclusions are discrete
comparisons of two points, not a continuous design response. (iii)
\emph{Transport radii.} Hard distance limits (150/250/400~km) approximate
continuous logistics costs; they follow literature cost benchmarks but truncate
the feasible region. (iv) \emph{Feedstock availability.} BT23 tonnages are
treated as fully collectable; field-level collection losses and residue
retention requirements are not modeled. A $-40\%$ availability sensitivity
costs 29\% of near-term surplus (the preferred corn-stover mix is
supply-tight) but only 4\% in the mature scenario, so the assumption matters
mostly for near-term deployment.
(v) \emph{GHG parameters.} The decomposition fraction (0.9), the permanence
factors (0.65/0.80, now aligned with VM0044 v1.2 Table 3), the
literature-based char H/C values (Table~S2; the eligibility outcome is inert
to $\pm$0.1 H/C and the first 500$^\circ$C flip needs $+0.25$, SI Section on
eligibility sensitivity), and the N$_2$O
meta-analysis mean (52~kg/t) are literature defaults; the decomposition
fraction is flagged for local calibration, and a sensitivity halving it to
0.5 lowers the \$200/t Paradigm~A surplus by 31--35\%
(Section~\ref{sec:sensitivity}); the
qualitative response and the paradigm ordering are robust, but the surplus
level is not, and A's magnitude advantage over the gated VM0044 credit narrows
as the decomposition fraction falls. (vi) \emph{VM0044 proxy.} Eligibility is
operationalized through the molar H/C gate only; additionality, leakage, and
monitoring/verification costs are outside the model scope. (vii)
\emph{Demand.} Demand segments are exogenous step functions calibrated to
market evidence rather than estimated elasticities, and their spatial
allocation across counties is proportional to farmland area rather than
observed demand locations. The binding demand constraint is therefore the
aggregate quantity, and the equilibrium facility pattern is feedstock-driven;
a demand-concentration sensitivity (SI, demand data pack) leaves the aggregate
conclusions unchanged. (viii) \emph{Static
structure.} The model is single-year with no seasonal inventory; the
allocation-design result (C3) and facility-path dependencies point to a
multi-period extension. (ix) \emph{Solver gaps.} Policy sweeps are reported at
0.5--4.7\% MIP optimality gaps (Table~S3). Every point whose original gap
exceeded 2\% was re-solved at a 0.1\% target with 3{,}600--7{,}200~s limits
(the near-term C2 $p_c=\$150$ point improved from a 4.7\% to a 2.5\% gap,
surplus $+1.9\%$; the v2b points improved to 1.4--2.7\%), forward versus
reverse sweep orders reproduce every point within 0.2\%, and the v2b Step-1
gaps are now recorded (2.6\% near-term, 1.0\% mature-market). The reported
structural responses are therefore stable across warm-start paths; TIME\_LIMIT
points remain feasible incumbents and their per-point gap, runtime, and status
are archived in the certificate CSVs. (x) \emph{Scope.} Wisconsin-specific magnitudes do not transfer
verbatim to other Midwest states (feedstock mix, demand, and dairy waste
streams differ); the qualitative policy mechanisms are the transferable
findings.
